% ZeroPP manuscript.
%
% Document class: plain article for now. Swap for the target journal's
% official LaTeX template at submission time (primary target: Machine
% Learning: Earth, IOP Publishing -- check iopart.cls / IOP's current
% author-template package once accepted for review; secondary targets QJRMS
% (Wiley) / AIES (AMS) have their own templates too. Do not submit with this
% generic class.
\documentclass[11pt]{article}

\usepackage[a4paper,margin=2.5cm]{geometry}
\usepackage[utf8]{inputenc}
\usepackage[T1]{fontenc}
\usepackage{lmodern}
\usepackage{amsmath,amssymb}
\usepackage{graphicx}
\usepackage{booktabs}
\usepackage{natbib}
\usepackage[breaklinks=true]{hyperref}
\usepackage{xcolor}

% Auto-generated numbers -- every figure/number cited in the text comes from
% here (scripts/generate_numbers.py reads results/ files and configs). Never
% hand-type a number from a results file directly in a .tex source file.
% AUTO-GENERATED by scripts/generate_numbers.py -- DO NOT EDIT BY HAND.
% Every value here is read from a real results/ file or config, never typed.
% Regenerate after any results/ file changes: python scripts/generate_numbers.py

\newcommand{\EuppbenchVersion}{v1.0}
\newcommand{\NStations}{49}
\newcommand{\NLeadTimes}{21}
\newcommand{\MaxLeadHours}{120}
\newcommand{\LeadStepHours}{6}
\newcommand{\NMatchedInstances}{737,809}
\newcommand{\QuantileLevelsList}{0.1, 0.2, 0.3, 0.4, 0.5, 0.6, 0.7, 0.8, 0.9}
\newcommand{\NQuantileLevels}{9}
\newcommand{\DaysPerCase}{3.4833}
\newcommand{\ReforecastTotalIssueDates}{209}
\newcommand{\ReforecastIssueDateSpanDays}{728}
\newcommand{\FullArchiveCases}{4180}
\newcommand{\DataSizeSweepSeeds}{5}
\newcommand{\TrainEnsembleMembers}{11}
\newcommand{\TestEnsembleMembers}{51}
\newcommand{\TrainYearRange}{1997--2016}
\newcommand{\TestYearRange}{2017--2018}
\newcommand{\NVarianceInflationVariants}{4}
\newcommand{\CrpsRawFull}{1.3590}
\newcommand{\CoverageRawFull}{0.4259}
\newcommand{\WidthRawFull}{2.3716}
\newcommand{\CrpsTsfmFull}{1.3592}
\newcommand{\CoverageTsfmFull}{0.7603}
\newcommand{\WidthTsfmFull}{5.1379}
\newcommand{\CrpsInflAFull}{1.3005}
\newcommand{\CoverageInflAFull}{0.7388}
\newcommand{\WidthInflAFull}{5.2037}
\newcommand{\CrpsInflBFull}{1.3042}
\newcommand{\CoverageInflBFull}{0.6096}
\newcommand{\WidthInflBFull}{3.7258}
\newcommand{\CrpsInflCFull}{1.3361}
\newcommand{\CoverageInflCFull}{0.8221}
\newcommand{\WidthInflCFull}{6.7224}
\newcommand{\CrpsEmosPooledFull}{1.2612}
\newcommand{\CoverageEmosPooledFull}{0.8303}
\newcommand{\WidthEmosPooledFull}{5.4559}
\newcommand{\CrpsEmosLocalFull}{1.1023}
\newcommand{\CoverageEmosLocalFull}{0.8289}
\newcommand{\WidthEmosLocalFull}{4.7548}
\newcommand{\CrpsDrnFull}{1.2251}
\newcommand{\CoverageDrnFull}{0.7720}
\newcommand{\WidthDrnFull}{4.6494}
\newcommand{\TwcrpsRawFull}{1.1109}
\newcommand{\TwcrpsTsfmFull}{1.0268}
\newcommand{\TwcrpsInflAFull}{0.9522}
\newcommand{\TwcrpsInflBFull}{0.9776}
\newcommand{\TwcrpsInflCFull}{1.0029}
\newcommand{\TwcrpsEmosPooledFull}{0.9138}
\newcommand{\TwcrpsEmosLocalFull}{0.8239}
\newcommand{\TwcrpsDrnFull}{0.8989}
\newcommand{\EmosPooledFitSeconds}{109.8}
\newcommand{\EmosLocalFitSeconds}{146.7}
\newcommand{\InflLeakyCrps}{1.3054}
\newcommand{\InflLeakyCoverage}{0.7603}
\newcommand{\InflLeakyWidth}{5.5344}
\newcommand{\InflLeakyWidthDiffPoint}{0.3965}
\newcommand{\InflLeakyWidthDiffCiLow}{0.2887}
\newcommand{\InflLeakyWidthDiffCiHigh}{0.4997}
\newcommand{\EnsSubsampleCrpsDiffPooled}{0.0121}
\newcommand{\EnsSubsampleWidthDiffPooled}{0.0131}
\newcommand{\EnsSubsampleCoverageDiffPooled}{0.0051}
\newcommand{\EnsSubsampleCrpsDiffLocal}{0.0141}
\newcommand{\EnsSubsampleWidthDiffLocal}{0.0171}
\newcommand{\EnsSubsampleCoverageDiffLocal}{0.0078}
\newcommand{\CrpsTsfmTwoCov}{1.3837}
\newcommand{\CoverageTsfmTwoCov}{0.8027}
\newcommand{\WidthTsfmTwoCov}{5.8246}
\newcommand{\CrpsTwoCovDiff}{0.0246}
\newcommand{\WidthTwoCovDiff}{0.6867}
\newcommand{\LeakyWidthAtTwoCovCoverage}{6.3026}
\newcommand{\LeakyWidthAtTwoCovCoverageDiff}{0.4780}
\newcommand{\SeasonalCompareNDays}{30}
\newcommand{\CoverageContigSeasonalN}{0.7668}
\newcommand{\CoverageRandomSeasonalN}{0.8105}
\newcommand{\CrpsContigSeasonalN}{1.2720}
\newcommand{\CrpsRandomSeasonalN}{1.2760}
\newcommand{\BpCrpsPooledCases}{1.6}
\newcommand{\BpCrpsPooledDays}{6}
\newcommand{\BpCrpsLocalCases}{1.8}
\newcommand{\BpCrpsLocalDays}{6}
\newcommand{\BpCoveragePooledCases}{8.4}
\newcommand{\BpCoveragePooledDays}{29}
\newcommand{\BpCoverageLocalCases}{56.2}
\newcommand{\BpCoverageLocalDays}{196}
\newcommand{\CoverageCrpsRatioPooled}{5.1}
\newcommand{\CoverageCrpsRatioLocal}{32.0}
\newcommand{\BpCrpsDrnCases}{15.1}
\newcommand{\BpCrpsDrnDays}{52}
\newcommand{\BpCoverageDrnCases}{181.4}
\newcommand{\BpCoverageDrnDays}{632}
\newcommand{\InflASmallestNDays}{30}
\newcommand{\InflASmallestNCrps}{1.2972}
\newcommand{\FirstSignFlipHours}{22.36}
\newcommand{\DurableCrossoverHours}{43.49}
\newcommand{\LeadZeroExactMatchFraction}{0.000000}
\newcommand{\LeadZeroMaxOtherExactMatchFraction}{0.000028}
\newcommand{\LeadZeroActualStd}{1.7644}
\newcommand{\LeadZeroPredictedStd}{1.7387}
\newcommand{\LeadZeroStdRatio}{1.0148}
\newcommand{\SigTestNDays}{30}
\newcommand{\SigTestNCases}{9}
\newcommand{\SigCrpsStationT}{0.1485}
\newcommand{\SigCrpsStationW}{1.88e-05}
\newcommand{\SigCrpsDayT}{1.50e-12}
\newcommand{\SigCrpsDayW}{3.59e-06}
\newcommand{\SigCoverageStationT}{0.7432}
\newcommand{\SigCoverageStationW}{0.0165}
\newcommand{\SigCoverageDayT}{0.0506}
\newcommand{\SigCoverageDayW}{0.4391}
\newcommand{\BpCrpsPooledSigK}{3}
\newcommand{\BpCrpsPooledSigDays}{10}
\newcommand{\BpCrpsPooledSigKDayT}{5.99e-07}
\newcommand{\BpCrpsPooledSigKDayW}{0.0186}
\newcommand{\BpCrpsPooledSigKStationT}{0.3446}
\newcommand{\BpCrpsPooledSigKStationW}{6.81e-05}
\newcommand{\BpCrpsPooledStationSigExists}{no}
\newcommand{\BpCrpsPooledStationWilcoxonFirstSigK}{2}
\newcommand{\BpCrpsLocalSigK}{3}
\newcommand{\BpCrpsLocalSigDays}{10}
\newcommand{\BpCrpsLocalSigKDayT}{7.51e-13}
\newcommand{\BpCrpsLocalSigKDayW}{1.03e-09}
\newcommand{\BpCrpsLocalSigKStationT}{0.0529}
\newcommand{\BpCrpsLocalSigKStationW}{3.49e-08}
\newcommand{\BpCrpsLocalStationSigExists}{yes}
\newcommand{\BpCrpsLocalStationWilcoxonFirstSigK}{2}
\newcommand{\LeadSigCrpsBucketShortDayDiff}{0.1682}
\newcommand{\LeadSigCrpsBucketShortDayT}{4.29e-61}
\newcommand{\LeadSigCrpsBucketShortDayW}{1.39e-63}
\newcommand{\LeadSigCrpsBucketShortStationT}{0.0354}
\newcommand{\LeadSigCoverageBucketShortDayDiff}{0.0573}
\newcommand{\LeadSigCoverageBucketShortDayT}{8.40e-53}
\newcommand{\LeadSigCoverageBucketShortDayW}{9.25e-49}
\newcommand{\LeadSigCoverageBucketShortStationT}{0.0255}
\newcommand{\LeadSigCrpsBucketMidDayDiff}{-0.0647}
\newcommand{\LeadSigCrpsBucketMidDayT}{4.33e-07}
\newcommand{\LeadSigCrpsBucketMidDayW}{0.0240}
\newcommand{\LeadSigCrpsBucketMidStationT}{0.2268}
\newcommand{\LeadSigCoverageBucketMidDayDiff}{0.0610}
\newcommand{\LeadSigCoverageBucketMidDayT}{6.15e-50}
\newcommand{\LeadSigCoverageBucketMidDayW}{1.96e-46}
\newcommand{\LeadSigCoverageBucketMidStationT}{0.0029}
\newcommand{\LeadSigCrpsBucketLongDayDiff}{-0.2453}
\newcommand{\LeadSigCrpsBucketLongDayT}{3.69e-35}
\newcommand{\LeadSigCrpsBucketLongDayW}{5.63e-34}
\newcommand{\LeadSigCrpsBucketLongStationT}{3.31e-08}
\newcommand{\LeadSigCoverageBucketLongDayDiff}{0.0837}
\newcommand{\LeadSigCoverageBucketLongDayT}{3.06e-64}
\newcommand{\LeadSigCoverageBucketLongDayW}{4.03e-58}
\newcommand{\LeadSigCoverageBucketLongStationT}{3.30e-06}
\newcommand{\CrpsPlainEmosMatched}{1.1362}
\newcommand{\CrpsPersistEmosMatched}{1.1312}
\newcommand{\CrpsTsfmMatchedShortLead}{0.9898}
\newcommand{\PersistenceWeight}{0.0786}
\newcommand{\PersistenceGapClosedPct}{3.4}
\newcommand{\NPersistenceMatchedInstances}{140,457}


\title{Zero-shot time series foundation models for statistical
  postprocessing of weather forecasts: how much training data do they
  replace?}
% Title decision (P3, panel review, 2026-09-07): Option A from
% docs/title_options.md, CONFIRMED FINAL by the author -- the plural
% "foundation models" stays, and the single-model-scope caveat this choice
% requires lives in the Introduction's closing paragraph and the
% Conclusion's closing paragraph instead of a title change. Do not revisit.

\author{\"{O}mer Faruk Avc{\i}\thanks{\texttt{avcio20@itu.edu.tr} \\
ORCID: \href{https://orcid.org/0009-0008-6942-7113}{0009-0008-6942-7113}}\\[4pt]
\parbox{0.85\linewidth}{\centering\small Department of Climate Science
and Meteorological Engineering, Faculty of Aeronautics and Astronautics,
\.{I}stanbul Technical University, \.{I}stanbul, T\"{u}rkiye}}
\date{}

\begin{document}
\maketitle

\begin{abstract}
Statistical postprocessing of ensemble forecasts is usually trained on
multi-year, per-station archives; a zero-shot time series foundation
model needs none. Using EUPPBench \EuppbenchVersion{}'s German 2\,m
temperature subset (\NStations{} stations, 0--\MaxLeadHours{}~h; all
methods scored on \NMatchedInstances{} matched instances), we
compare a frozen TimesFM-3 (ensemble mean as a past-future covariate)
against Ensemble Model Output Statistics (EMOS) and a one-parameter
variance-inflation rescaling. At full training data, every trained
method beats TimesFM-3 on the Continuous Ranked Probability Score
(CRPS); so does a fixed-climatological rescaling using no training data
at all. Pooled EMOS's CRPS crosses TimesFM-3's at \BpCrpsPooledCases{}
cases (\BpCrpsPooledDays{} days), though the crossing is significant
only from \BpCrpsPooledSigK{} cases (\BpCrpsPooledSigDays{} days)
onward, for both EMOS variants, under day-blocking.
80\%-nominal coverage crosses more slowly (\CoverageCrpsRatioPooled{}x
and \CoverageCrpsRatioLocal{}x later than CRPS, pooled and local EMOS).
TimesFM-3 retains a calibration-sharpness advantage: even a
leakage-optimistic multiplier fit on test-set coverage needs
\InflLeakyWidthDiffPoint{} K more at matched coverage (95\% CI
[\InflLeakyWidthDiffCiLow{}, \InflLeakyWidthDiffCiHigh{}]). Its
short-lead CRPS advantage reverses by \DurableCrossoverHours{}~h; a
persistence-augmented EMOS closes only \PersistenceGapClosedPct{}\% of
the gap, which does not support an information-asymmetry
explanation, though not from architecture. Supplying ensemble spread as
a second covariate does not close the calibration gap either: it
worsens CRPS and reaches nominal coverage only with a wider interval.
Overall, on CRPS this
frozen model is mostly not competitive with a few days of
station-specific training data -- with one exception that survives even
a leakage-optimistic bound, and a short-lead advantage that the
information-asymmetry explanation does not account for.
\end{abstract}

\section{Introduction}
\label{sec:intro}

Numerical weather prediction (NWP) output is systematically biased at any
single station or grid point: model resolution, physical
parameterizations, and terrain representation err in a station-specific
and often persistent way, and statistical postprocessing exists
specifically to correct this systematic component rather than relying on
raw ensemble output directly \citep{vannitsem_2020_postprocessing_review,gneiting_2005_emos}.
Beyond Ensemble Model Output Statistics (EMOS)'s original Gaussian
regression \citep{gneiting_2005_emos}, the field has extended
postprocessing along several axes: mixture and non-Gaussian predictive
families \citep{baran_2015_mixture_emos}, target-specific reformulations
for variables EMOS's Gaussian assumption fits poorly
\citep{scheuerer_2013_precip_emos}, explicit calibration-enforcing
extensions \citep{wilks_2018_enforcing_calibration}, and
permutation-invariant neural architectures that treat ensemble members as
an unordered set \citep{hoehlein_2023_permutation_invariant}. The
Distributional Regression Network (DRN, \citealp{rasp_2018_drn}) is this
paper's own neural baseline (Section~\ref{sec:methods}).

Every one of these methods shares the same requirement: a multi-year
archive of paired ensemble/observation history at the specific station
being corrected. \citet{vannitsem_2020_postprocessing_review}'s review of
the field flags this as a structural limitation, not an incidental one --
and it is exactly what a newly commissioned station, a station following
an NWP model version upgrade that breaks the archive's stationarity, or a
station in a sparsely observed region cannot supply. Trained
postprocessing's accuracy gains, in other words, come with an entry cost
that a purely zero-shot alternative would not.

Time series foundation models (TSFMs) -- large, pretrained sequence models
applied without any task-specific training or parameter updates on the
target series -- offer a structurally different answer: no
station-specific training archive at all. This is a different proposition
from the spatial, full-field AI weather-prediction models that have
separately transformed NWP itself
\citep{lam_2023_graphcast,price_2024_gencast,bodnar_2025_aurora}: those
models replace the numerical forecast; a TSFM used for postprocessing
takes an existing ensemble forecast as input and corrects it at the
station level, one series at a time, with no spatial field anywhere in its
architecture. TimesFM-3 -- the 3.0 release of Das et al.'s decoder-only
TimesFM architecture \citep{das_2024_timesfm_icml} -- is one such model, and
its past-future covariate support (a version-specific capability sourced
directly in Section~\ref{sec:covariates}, not part of the original
architecture paper) is the specific capability this paper exploits: the
ensemble forecast can be supplied as a covariate known across the entire
forecast horizon, not just in the past, letting a zero-shot model use NWP
guidance the way a postprocessing method would, without training on this
benchmark's data at all.

To our knowledge, no existing work directly measures the training-data
size at which a station's own trained postprocessing model overtakes a
zero-shot time series foundation model given the same NWP guidance -- the
question this paper answers. On the identical EUPPBench benchmark this
paper uses, \citet{vanpoecke_2024_transformer},
\citet{mlakar_2023_spline_flows}, and \citet{landry_2025_fmap} each
propose a trained neural or flow-based postprocessing architecture, and
\citet{allen_2023_multivariate_calibration} evaluates multivariate
calibration on it -- none compares against a zero-shot foundation model.
\citet{baran_2026_sharpness_penalty} studies the same sharpness axis this
paper's calibration-sharpness result depends on, on the same benchmark,
but not the zero-shot-vs-trained breakpoint question.
\citet{sun_2026_streamflow_zeroshot} pose a structurally identical
question -- is a zero-shot time series foundation model competitive with
domain-trained methods -- in hydrology, not weather postprocessing;
\citet{bharadwaj_2026_air_quality_arena} benchmark TSFM zero-shot skill
across regions in air quality forecasting, a related but distinct target
variable and evaluation design. \citet{wessel_2024_leadtime_continuous}
and \citet{muschinski_2023_mos_random_forests} are the closest
small-training-data postprocessing rivals to this paper's own breakpoint
curve, targeting the same low-data regime with different
(lead-time-continuous, random-forest) trained architectures rather than a
zero-shot comparison. None of these measures where the crossing point
actually is.

This paper's contribution is threefold. (1) At full training data, every
trained method (EMOS pooled, EMOS local, DRN) beats TimesFM-3's
Continuous Ranked Probability Score (CRPS, Section~\ref{sec:results-full}).
A one-parameter variance-inflation
rescaling of the raw ensemble does too: a fixed-climatological variant
beats it with no training data at all, and the CRPS-optimal variant beats
it at every training size we tested it at, from \InflASmallestNDays{}
days upward (Section~\ref{sec:results-full}) -- on CRPS, the answer to
this paper's central question is negative for TimesFM-3, and not
narrowly so. This advantage starts at very small training sizes: pooled
EMOS's CRPS crosses TimesFM-3's at \BpCrpsPooledCases{} cases
(\BpCrpsPooledDays{} calendar days, Section~\ref{sec:results-breakpoint}),
but that raw crossing is not itself statistically significant; the
advantage is significant from \BpCrpsPooledSigK{} cases
(\BpCrpsPooledSigDays{} calendar days) onward under day-blocking, our
primary block definition, for both EMOS pooled and EMOS local alike --
the station-blocked \(t\)-test does not reach significance anywhere in
the tested \(k=1\)--9 range, while the station-blocked Wilcoxon test does
from \(k=\BpCrpsPooledStationWilcoxonFirstSigK{}\) onward, the same
test-statistic divergence already reported at \SigTestNCases{} cases
(Section~\ref{sec:results-significance}), now shown to persist across the
whole low-\(k\) range. Coverage crosses more slowly
(ratios computed from the unrounded breakpoint estimates, not the
values as displayed): pooled EMOS's coverage crosses TimesFM-3's roughly
\CoverageCrpsRatioPooled{}x later than its CRPS does (\BpCoveragePooledCases{}
cases / \BpCoveragePooledDays{} days), and local EMOS's about
\CoverageCrpsRatioLocal{}x later (\BpCoverageLocalCases{} cases /
\BpCoverageLocalDays{} days) -- so "a few days is enough" holds for CRPS
specifically, not for calibration. These comparisons are between methods
given different information: TimesFM-3's covariate channel received the
ensemble mean but not ensemble spread in this study, while EMOS and DRN
received spread but no observation history (Section~\ref{sec:covariates},
Table~\ref{tab:info-channels}); supplying spread as a second covariate
does not close this calibration-sharpness gap either
(Section~\ref{sec:discussion-covariate-check}). (2) On the
calibration-sharpness axis, TimesFM-3 holds a real advantage over a
trivial rescaling: even a leakage-optimistic multiplier fit directly on
test-set coverage -- an upper bound no legitimately deployable trivial
rescaling could exceed -- needs \InflLeakyWidth{} K, still wider than
TimesFM-3's \WidthTsfmFull{} K (Section~\ref{sec:discussion-trivial}).
(3) TimesFM-3's short-lead CRPS advantage first
flips sign at \FirstSignFlipHours{}~h and reverses durably by
\DurableCrossoverHours{}~h (Section~\ref{sec:results-leadtime}); the
mechanism behind that advantage has two independent parts. The
observation-history part was tested directly and does not support an
information-asymmetry explanation -- a persistence-augmented EMOS
closes only \PersistenceGapClosedPct{}\% of the gap -- but the test
itself cannot cleanly distinguish an information-set explanation from an
architectural one (Section~\ref{sec:discussion-mechanism}). The
ensemble-spread/calibration part was also tested directly, by supplying
TimesFM-3 with ensemble spread as a second covariate: it worsens CRPS
and shifts the model along the calibration-sharpness trade-off without
closing the calibration deficit, even though the channel demonstrably
accepts and uses the input (Section~\ref{sec:discussion-covariate-check}).
Together these results argue that a frozen, zero-shot time series
foundation model is, for this task, mostly not competitive with a few
days of station-specific training data on CRPS -- with one genuine
exception that survives even a leakage-optimistic trivial baseline (the
calibration-sharpness axis), and a short-lead advantage that the
information-asymmetry explanation does not account for. The evidence
presented here is for one
time series foundation model, TimesFM-3; whether it generalizes to other
pretrained time series models (e.g.\ Chronos-2, Moirai-2) is untested.

\section{Data and experimental design}
\label{sec:data}

We use the EUPPBench postprocessing benchmark, version~\EuppbenchVersion{}
\citep{demaeyer_2023_euppbench}, restricted to its German subset of
\NStations{} stations and the \(2\,\mathrm{m}\) temperature target. EUPPBench
defines its own train/test split, which we never re-partition: training
uses ECMWF reforecasts issued \TrainYearRange{} with
\TrainEnsembleMembers{} ensemble members, and testing uses operational
ECMWF forecasts issued \TestYearRange{} with \TestEnsembleMembers{}
members. Both sides carry \NLeadTimes{} lead times from 0 to
\MaxLeadHours{}~h in \LeadStepHours{}-hour steps. Reference-quality
observations are DWD station measurements; the ensemble side is ECMWF's
operational and reforecast archives as distributed through EUPPBench.
EUPPBench also defines a wind-speed target (\(w10\)); we defer it here
because no country in the benchmark provides a direct wind-speed
observation series, only wind gust, and substituting gust for speed would
change what is being verified.

TimesFM-3's fixed context-window requirement (at least 40 prior
observations) excludes some station/lead/issue-time combinations that would
otherwise be evaluable by the other methods, so every comparison in this
paper is computed on the intersection of instances all methods can produce
a prediction for, not each method's own largest possible set. The full test
set otherwise available to EMOS and the raw ensemble contains 748{,}099
instances; requiring TimesFM-3's context window reduces this to
\NMatchedInstances{} matched instances, and every result in this paper is
computed on that matched set for every method, including EMOS and the raw
ensemble, not their own larger native set.

The archive's raw station dimension carries 51 station identifiers; two are
excluded throughout this study because their observation series contain
missing values that would otherwise propagate NaNs into every downstream
metric, leaving the \NStations{} stations used everywhere in this paper.
Ensemble summary statistics (mean and variance) are computed once from the
raw per-member forecasts before any postprocessing method sees them, using
an unbiased (\(\mathrm{ddof}=1\)) variance estimator on both the training
and test sides so that EMOS is never trained against a systematically
different variance definition than the one it is evaluated against. The
training-side ensemble variance is nonetheless estimated from
\TrainEnsembleMembers{} members against \TestEnsembleMembers{} on the test
side; an 11-member variance estimate is intrinsically noisier than a
51-member one of the same true spread, so EMOS is trained against a
noisier estimate of ensemble variance than the one used to evaluate it
(Section~\ref{sec:limitations}).

The paper's central axis is training-data size, expressed as a number of
training \emph{cases} \(k\) rather than calendar days directly, since the
reforecast archive samples a fixed, non-daily grid of issue dates rather
than every calendar day. The reforecast archive's raw \texttt{time} coordinate cannot be reliably
decoded to calendar dates: attempting to decode it raises an
\texttt{OverflowError}, confirmed to be caused by an unmasked fill-value
sentinel surviving into the time-unit conversion, not by a transient
network fault, since the same failure text was independently confirmed on
the sibling reforecast-observations file. We therefore treat each issue
date as an abstract
\((\mathrm{year\_idx}, \mathrm{time\_idx})\) pair rather than a decoded
calendar date: \ReforecastTotalIssueDates{} representative issue-date
positions (\(\mathrm{time\_idx}\)) recur once per analog year
(\(\mathrm{year\_idx}\), 20 analog years in total), so the full archive
contains \FullArchiveCases{} cases (\(\ReforecastTotalIssueDates \times
20\)). We measured the calendar-day span of one such 209-date cycle
directly from its real, decodable date labels: it spans
\ReforecastIssueDateSpanDays{} calendar days --
consistent with a twice-weekly issue cadence over two calendar years, not
one -- giving \(\DaysPerCase\) calendar days per case. We convert a
requested training size of \(N\) days to \(k = \mathrm{round}(N /
\DaysPerCase)\) cases throughout, and report both units together wherever
a training size is cited, since case count and calendar-day equivalent are
not interchangeable at this archive's actual sampling density. The
contiguous training arm (below) orders cases by descending
\((\mathrm{year\_idx}, \mathrm{time\_idx})\), assuming ascending
\(\mathrm{year\_idx}\) is chronologically more recent; only
\(\mathrm{time\_idx}\)'s ordering was independently verified against real
decoded dates, so this assumption is disclosed, not confirmed
(Section~\ref{sec:limitations} bounds its limited practical impact).

We test two ways of selecting which \(k\) cases make up a given training
size, since they answer different questions. The \emph{contiguous} arm
(primary) takes the \(k\) cases closest in time to the test period,
working backward -- the realistic scenario for a newly deployed station
whose limited history is chronologically contiguous, so small \(k\) sees
substantially less seasonal variety than a full archive would. The
\emph{random} arm (secondary, \DataSizeSweepSeeds{} seeds, mean and
standard deviation reported) draws \(k\) cases uniformly at random from
the full archive, giving artificial full-season coverage even at small
\(k\). The gap between the two arms' results is itself a reported finding
about the operational value of seasonal training coverage
(Section~\ref{sec:results}), not a discrepancy to be resolved in favor of
one arm. \citet{baran_2026_sharpness_penalty} targets the same sharpness
axis on the same benchmark; we discuss the relationship to our own
sharpness results directly in Section~\ref{sec:discussion} rather than
here. Other EUPPBench applications are surveyed in
Section~\ref{sec:intro}.

\section{Methods}
\label{sec:methods}

\subsection{Postprocessing methods compared}
\label{sec:methods-list}

We compare six families of methods, all sharing a single interface: given a
station's ensemble forecast summary for a lead time, produce a predictive
distribution for the observed \(2\,\mathrm{m}\) temperature at that lead,
expressed as the \NQuantileLevels{} quantiles \QuantileLevelsList{} (a fixed
global constant used identically by every method in this study).

\textbf{Raw ensemble.} The empirical distribution of the raw NWP ensemble
members at each station and lead time, used unmodified. This is the
zero-shot, zero-parameter baseline every postprocessing method is expected to
improve on.

\textbf{EMOS.} Ensemble Model Output Statistics \citep{gneiting_2005_emos}, a
nonhomogeneous Gaussian regression fit by minimizing the closed-form Gaussian
CRPS \citep{gneiting_2007_proper_scoring} over the training set:
\(\mu = a + b \cdot \bar{x}\), \(\sigma^2 = \exp(c) + \exp(d) \cdot s^2\),
where \(\bar{x}\) and \(s^2\) are the ensemble mean and variance. We fit two
variants of this same model, differing only in how many stations share one
parameter set: \emph{pooled} EMOS fits one global \((a,b,c,d)\) across all
\NStations{} stations, following the field's standard practice
\citep{gneiting_2005_emos,raftery_2005_bma}; \emph{local} EMOS fits a
separate \((a,b,c,d)\) per station whenever that station has enough training
rows to do so (we require at least 5 rows; stations below this threshold are
excluded from local EMOS's evaluation at that training size, and the
resulting station-coverage fraction is reported alongside every local-EMOS
result). \citet{baran_2016_combining} and \citet{gebetsberger_2018_estimation}
report comparable station-adaptive EMOS variants at small training sizes,
the regime this paper's breakpoint curve targets directly.

\textbf{DRN.} The Distributional Regression Network of
\citet{rasp_2018_drn}: a feed-forward neural network with a learned
per-station embedding, taking the ensemble mean and variance as input and
outputting \((\mu, \log\sigma^2)\) parameters of the same Gaussian family
EMOS uses, trained by gradient descent (Adam) on the same closed-form
Gaussian CRPS loss, reimplemented in PyTorch so it is differentiable. Unlike
the original DRN, which additionally uses auxiliary NWP fields and lead time
as inputs, our training pipeline pools all \NLeadTimes{} lead times into one
fit and does not expose lead time to any model (Section~\ref{sec:limitations});
DRN and EMOS are therefore compared under the same feature constraint, not
DRN's originally-intended full feature set. \citet{schulz_2021_windgusts}
runs a systematic comparison of neural postprocessing architectures under
this same feature-constrained setup, the closest methodological precedent
to our own DRN evaluation. Time-series/lead-time-continuous EMOS variants
\citep{jobst_2024_ts_emos,wessel_2024_leadtime_continuous} target the
small-training-data and lead-time-resolved regimes this paper's breakpoint
curve also explores, without themselves being baselines here (a wider
survey of the neural-postprocessing and quantile-regression literature is
given in Section~\ref{sec:intro}).

\textbf{TimesFM-3 (zero-shot).} The 3.0 release of TimesFM, a frozen,
pretrained decoder-only time series foundation model. The architecture
family is described in \citet{das_2024_timesfm_icml}; that paper describes the
original TimesFM 1.0 decoder-only architecture and does not cover the
3.0 release's covariate support, which Section~\ref{sec:covariates}
sources directly rather than through that citation. This project uses
checkpoint \texttt{google/timesfm-3.0-pytorch}, with no fine-tuning,
task-specific training, or parameter updates of any kind.
Section~\ref{sec:covariates} describes its input construction. Contrasted
with generalist AI weather-prediction models that operate on full spatial
fields \citep{lam_2023_graphcast,price_2024_gencast,bodnar_2025_aurora} --
TimesFM-3 is used here in univariate, single-station mode, not as a
spatial forecasting system, and is not compared against those models
here.

\textbf{Variance-inflation baseline.} A trivial, zero/low-parameter
alternative: scale the raw ensemble's per-instance spread by a single
multiplier \(\lambda\) and read off Gaussian quantiles at the same
\NQuantileLevels{} levels. We report \NVarianceInflationVariants{} variants,
each answering a different question about whether TimesFM-3's apparent skill
is trivially reproducible \citep{zhang_2025_context_parroting}: (a) a
CRPS-optimal \(\lambda\) fit by minimizing training-set CRPS at each training
size \(N\) (fair comparison against EMOS, which also uses \(N\) training
cases); (b) a fixed climatological \(\lambda\), requiring no data at all
(fair comparison against TimesFM-3's own zero-shot status); (c) a
coverage-matched \(\lambda\), fit on training data only (never on test-set
coverage) by root-finding to reach as closely as a training-fit multiplier
allows TimesFM-3's own empirical coverage, then evaluated for interval
width on the test set -- this directly asks whether TimesFM-3's
calibration, not just its CRPS, is achievable by a one-parameter
rescaling. Because \(\lambda\) is fit on training data and applied to a
different test set, the coverage this variant actually reaches on test
does not land exactly on the training-time target; we report both values
together (Section~\ref{sec:results-full}) rather than the word
``exactly.'' (d) a leakage-optimistic \(\lambda\), fit by root-finding
directly on \emph{test-set} coverage (never on training data) to reach
TimesFM-3's own empirical coverage -- the deliberately
leakage-contaminated counterpart to (c), answering not "what can a
legitimate, leakage-free trivial rescaling achieve" but "what is the
ceiling on what any trivial rescaling could ever achieve at this
coverage target, even with an unfair advantage." Variant (d) is never
presented as a legitimate baseline result; it exists solely to bound how
much of variant (c)'s width disadvantage against TimesFM-3
(Section~\ref{sec:results-full}) could in principle be a leakage
artifact of the train/test split rather than a real limit of trivial
rescaling.

\subsection{TimesFM-3 covariate construction, and the information asymmetry it creates}
\label{sec:covariates}

TimesFM-3's target series is a station's own observation history; its
past-future (dynamic) covariate channel\footnote{Native past-future
covariate support is documented for the 3.0 release specifically, not in
the original TimesFM architecture paper \citep{das_2024_timesfm_icml}: model
card at \url{https://huggingface.co/google/timesfm-3.0-pytorch} and
release notes at
\url{https://research.google/blog/timesfm-3-a-zero-shot-foundation-model-for-multivariate-forecasting/}
(Google Research; both accessed 2026-09-07).} is filled with the ensemble mean,
supplied as a value known across the entire forecast horizon rather than
only in the past. This study supplied only the ensemble mean to this
channel; ensemble spread was not. This is an input-design choice, not a
confirmed architectural limit: the installed package's own source
documents the channel's accepted shape as genuinely multi-channel, and a
preliminary test confirms the model's output is measurably sensitive to
a second channel's real informational content, not merely its presence
(Section~\ref{sec:limitations} reports this preliminary test;
Section~\ref{sec:discussion-covariate-check} reports the full evaluation
it motivated). The model is entirely frozen: no fine-tuning,
no gradient updates, and no task-specific training of any kind are
performed at any point in this study.
\citet{kreusel_2026_covariates_key} independently reports that covariate
handling is the dominant factor in time series foundation model
performance; a coding error encountered during this project's own
development, where an incoherent covariate channel measurably degraded
TimesFM-3's apparent skill before being corrected, is consistent with that
finding (Section~\ref{sec:pitfall}).

This covariate construction leaves TimesFM-3 and the trained statistical
methods (EMOS, DRN) with genuinely different information, as this study
was actually run: TimesFM-3 sees the station's own recent observation
history but was not given ensemble spread; EMOS and DRN see ensemble
spread but no observation history at all (Table~\ref{tab:info-channels}).
We treat this asymmetry as a candidate explanation for three separate
empirical results in this paper (Section~\ref{sec:results}): TimesFM-3's
short-lead advantage (recent observations carry strong persistence
information at short lead times, information EMOS/DRN were not given
observation history to use in this study's configuration),
its calibration deficit relative to EMOS (TimesFM-3's covariate channel
was not given ensemble spread, the direct measure of forecast
uncertainty, in this study), and its long-lead disadvantage (observation
history loses predictive value as lead time grows, leaving TimesFM-3's
remaining information -- the ensemble mean alone -- no richer than what
EMOS already uses). This explanation was conditional on whether ensemble spread could be
added to TimesFM-3's covariate channel without degrading it -- a
question this project has since tested directly
(Section~\ref{sec:discussion-covariate-check}). The channel does accept
a second, simultaneous input (Section~\ref{sec:limitations}), but
supplying spread this way did not close the calibration deficit: CRPS
and interval width both worsen, and coverage moves only marginally
closer to the nominal target at a substantially wider interval. This
does not confirm the information-asymmetry framing for this half of the
mechanism question, but it does not confirm an architectural explanation
either (Section~\ref{sec:discussion-covariate-check},
Section~\ref{sec:conclusion}).
Section~\ref{sec:discussion} tests both halves of this explanation
directly: the observation-history half by augmenting EMOS with a
persistence predictor, and the calibration half by supplying TimesFM-3
with ensemble spread as a second covariate.

\begin{table}[htbp]
\centering
\caption{Information supplied to each method in this study -- not each
method's inherent architectural capability. TimesFM-3's ``no'' for
ensemble spread reflects this study's covariate design, not a confirmed
limit of the 3.0 release (see note below and Section~\ref{sec:limitations}).}
\label{tab:info-channels}
\begin{tabular}{lccc}
\toprule
Method & Observation history & Ensemble mean & Ensemble spread \\
\midrule
EMOS (pooled / local) & no & yes & yes \\
DRN & no & yes & yes \\
TimesFM-3 (zero-shot) & yes & yes & no$^{\ast}$ \\
Raw ensemble & no & yes (empirical) & yes (empirical) \\
Variance-inflation baseline & no & no & yes (scaled) \\
\bottomrule
\end{tabular}

\smallskip
{\small $^{\ast}$Not supplied in this study. TimesFM-3's covariate
parameter accepts a second, simultaneous channel by design (see the note
above), and a full evaluation supplying it found CRPS and interval width
both worsen, not improve (Section~\ref{sec:discussion-covariate-check}).}
\end{table}

Two other covariate-aware zero-shot approaches make a different design
choice than the frozen, built-in covariate channel this paper relies on:
\citet{auer_2025_icl_covariates} inject covariates through in-context
examples rather than an architectural slot, requiring no model changes
but also no guarantee the model attends to them the way a dedicated
covariate channel would; \citet{qin_2025_cora} instead adapt a pretrained
TSFM's covariate handling per-task through a small tuning step, which
this paper's zero-shot, no-fine-tuning design (Section~\ref{sec:methods-list})
deliberately avoids.

\subsection{Metrics}
\label{sec:metrics}

The Continuous Ranked Probability Score (CRPS) is our primary metric.
Every reported CRPS value in this paper, for every method, is computed by
the same pinball-loss-based quantile estimator
\citep{jordan_2017_scoringrules} applied to the identical
\NQuantileLevels{}-level quantile representation: EMOS and DRN's own
closed-form Gaussian CRPS \citep{gneiting_2007_proper_scoring,hersbach_2000_crps}
is used only inside their training objective (Section~\ref{sec:methods-list}),
never for a reported evaluation metric, so no method receives an
analytic-CRPS advantage over the discretized estimator every other method
is necessarily scored with. We additionally report the tail-weighted CRPS
(twCRPS), the same pinball-loss estimator restricted to the upper two
quantile levels (0.8, 0.9) \citep{allen_2023_weighted_scoringrules,taillardat_2019_extremes},
as a check on tail behavior; \citet{zamo_2018_crps_estimation},
\citet{bolin_2019_scale_invariance}, and \citet{arnold_2023_crps_decomposition}
discuss estimator choice and decomposition properties of CRPS more broadly,
and \citet{leutbecher_2020_crps_changes} interprets CRPS differences
directly. We additionally report \emph{coverage@80\%}: the empirical
fraction of observations falling within the predicted p10--p90 interval,
against a nominal level of 0.80 -- not 0.90, since our fixed
\NQuantileLevels{}-level quantile grid has no p05/p95 pair to compute a
90\%-nominal interval from; every mention of coverage in this paper refers
to this 80\%-nominal band and is labeled accordingly throughout. Finally we
report \emph{interval width} (the p10--p90 range in Kelvin) as a measure of
sharpness. Per \citet{gneiting_2007_calibration_sharpness}'s principle of
maximizing sharpness subject to calibration, interval width is reported as a
curve for every method and is never assigned a single ``better'' direction
or a breakpoint of its own (Section~\ref{sec:breakpoint-def}): a narrower
interval is only meaningful jointly with adequate coverage, not on its own.
PIT histograms (Section~\ref{sec:results}) are likewise reported
descriptively; \citet{broecker_2018_serial_dependence} and
\citet{broecker_2020_stratified_rank} show that formal rank-uniformity
hypothesis tests require correction for serial dependence, which our PIT
values genuinely exhibit, so no such test is applied here
(Section~\ref{sec:limitations}). \citet{dirkson_2025_misdiagnosing_reliability}
further shows that coverage and rank-based reliability diagnostics can both
be misled by the same covariance pathology; we do not claim these
descriptive diagnostics are immune to that critique (Section~\ref{sec:limitations}).

\subsection{Statistical significance}
\label{sec:significance}

The evaluation set is a station~$\times$~time~$\times$~lead panel, not a
single independent sample: instances share temporal autocorrelation within
a station and spatial (synoptic) dependence across the \NStations{}
stations on a given day. A textbook Diebold--Mariano test
\citep{diebold_1995_comparing_accuracy,diebold_2015_twenty_years} assumes one
autocorrelated series and is not directly applicable to this panel structure;
its small-sample behavior is separately unreliable even in the single-series
case it was designed for \citep{coroneo_2020_small_samples}. We test two
block definitions, aggregating the per-instance loss differential to one
mean value per block before testing: a \emph{station block} (absorbing
within-station temporal correlation, but not cross-station dependence) and
a \emph{day block} (absorbing cross-station synoptic dependence on a given
day, following the logic of exact tests under contemporaneous correlation,
\citealp{luger_2004_exact_tests}, but not multi-day persistence in the loss
differential itself). We do not additionally test a multi-day moving-block
variant that would absorb such short-range temporal persistence; the day
block already targets the specific mechanism this analysis was designed to
address (same-day synoptic systems affecting many stations at once), and we
report this as a remaining, open block-definition question rather than as a
claim that day-blocking is fully sufficient. For each block definition we
report a paired \(t\)-test and a Wilcoxon signed-rank test on the block
means, and a block-bootstrap confidence interval on the CRPS skill score
(resampling whole blocks with replacement). We call this a
\emph{station-blocked} or \emph{day-blocked paired test}, not a
Diebold--Mariano test: the classical test's assumptions do not transfer to
this station~$\times$~time~$\times$~lead panel.

\subsection{Breakpoint definition}
\label{sec:breakpoint-def}

The paper's central quantity is the training-data size at which a trained
method's performance crosses a zero-shot reference's. We define \(N\) days
of training as \(k = \mathrm{round}(N / \DaysPerCase)\) issue-date cases,
\DaysPerCase{} being the reforecast archive's measured ratio of calendar
days to sampled cases (Section~\ref{sec:data}). A breakpoint is
the linearly interpolated case count \(k\) at which a trained method's
metric crosses the zero-shot reference's value, computed over an ordered
sequence of tested \(k\); we distinguish three outcomes explicitly rather
than collapsing them into one ``no breakpoint'' label, which conflates
opposite situations: \emph{a crossing at \(k=X\)} (reported with its
calendar-day equivalent), \emph{already better at the smallest tested
\(k\)} (no crossing needed because the trained method starts ahead), and
\emph{no crossing within the tested range} (the trained method remains
behind throughout every tested \(k\), which is a distinct and, for the
zero-shot model's defense, more informative outcome than a numerical
breakpoint would be).

\section{Results}
\label{sec:results}

\subsection{Full-training-data comparison}
\label{sec:results-full}

% AUTO-GENERATED by scripts/generate_numbers.py -- DO NOT EDIT BY HAND.
\begin{table}[htbp]
\centering
\caption{Full-training-data (N=full, 4180 cases) comparison. CRPS, twCRPS, and interval width in Kelvin; coverage against the 80\%-nominal p10--p90 band. DRN is trained under a feature-constrained setup (Section~\ref{sec:limitations}), not its originally-intended full feature set.}
\label{tab:full-comparison}
\setlength{\tabcolsep}{4pt}
\begin{tabular}{lcccc}
\toprule
Method & \shortstack{CRPS\\(K)} & \shortstack{twCRPS\\(K)} & \shortstack{Coverage\\@80\%} & \shortstack{Interval\\width (K)} \\
\midrule
Raw ensemble & 1.3590 & 1.1109 & 0.4259 & 2.3716 \\
Var.-inflated, CRPS-optimal (a) & 1.3005 & 0.9522 & 0.7388 & 5.2037 \\
Var.-inflated, fixed $\lambda=1.5$ (b) & 1.3042 & 0.9776 & 0.6096 & 3.7258 \\
Var.-inflated, coverage-matched (c) & 1.3361 & 1.0029 & 0.8221 & 6.7224 \\
Var.-inflated, leakage-optimistic (d)$^{\dagger}$ & 1.3054 & 0.9579 & 0.7603 & 5.5344 \\
TimesFM-3 (zero-shot) & 1.3592 & 1.0268 & 0.7603 & 5.1379 \\
TimesFM-3, two-covariate (mean + spread)$^{\ddagger}$ & 1.3837 & -- & 0.8027 & 5.8246 \\
EMOS pooled & 1.2612 & 0.9138 & 0.8303 & 5.4559 \\
EMOS local & 1.1023 & 0.8239 & 0.8289 & 4.7548 \\
DRN & 1.2251 & 0.8989 & 0.7720 & 4.6494 \\
\bottomrule
\end{tabular}
\par\smallskip
{\small $^{\dagger}$Never a legitimate baseline: $\lambda$ is fit directly on test-set coverage (Section~\ref{sec:methods-list}), an upper bound on what any trivial rescaling could achieve at this coverage target, not a leakage-free result.}\\
{\small $^{\ddagger}$Not this paper's main TimesFM-3 configuration (that is the zero-shot row above, used everywhere else in this paper): an additional experiment supplying ensemble spread as a second covariate (Section~\ref{sec:discussion-covariate-check}). twCRPS is omitted (--) because per-instance quantiles for this run were not persisted; only the aggregate metrics shown were saved.}
\end{table}


At full training data (Table~\ref{tab:full-comparison}), every trained
method beats TimesFM-3 on CRPS: EMOS pooled (\CrpsEmosPooledFull{} K),
EMOS local (\CrpsEmosLocalFull{} K), and DRN (\CrpsDrnFull{} K), against
TimesFM-3's \CrpsTsfmFull{} K and the raw ensemble's \CrpsRawFull{} K.
The trivial variance-inflation baseline beats TimesFM-3 on CRPS too: both
the CRPS-optimal variant (a, \CrpsInflAFull{} K) and the
fixed-climatological variant (b, \CrpsInflBFull{} K) undercut TimesFM-3's
\CrpsTsfmFull{} K, and (b) needs no training data at all -- a single
zero-shot scalar rescaling of the raw ensemble's spread already beats a
frozen foundation model on this metric. On the calibration-sharpness axis
the same trivial rescaling does not match TimesFM-3: the coverage-matched
variant (c), calibrated on training data alone to reach as closely as a
training-fit multiplier allows TimesFM-3's own test coverage
(\CoverageInflCFull{} achieved against a \CoverageTsfmFull{} target,
Section~\ref{sec:methods}), needs \WidthInflCFull{} K to do so --
substantially wider than TimesFM-3's \WidthTsfmFull{} K, even after
accounting for (c) landing closer to the nominal 0.80 target than
TimesFM-3 itself, which alone would call for a somewhat wider interval.
Variant (d), never a legitimate baseline (Section~\ref{sec:methods-list}),
bounds how much of that width gap could be a leakage artifact of (c)'s
train/test split rather than a real limit: fit directly on test-set
coverage, it reaches \InflLeakyCrps{} K on CRPS and
\InflLeakyCoverage{} achieved against the identical \CoverageTsfmFull{}
target -- matched exactly, not merely approximately -- yet still needs
\InflLeakyWidth{} K, wider than TimesFM-3's \WidthTsfmFull{} K even at
this leakage-optimistic upper bound. This margin is not just a point
estimate: a day-blocked bootstrap on the per-instance width difference
gives \InflLeakyWidthDiffPoint{} K (95\% CI
[\InflLeakyWidthDiffCiLow{}, \InflLeakyWidthDiffCiHigh{}]), entirely
positive -- the leakage-optimistic variant is significantly wider than
TimesFM-3, not merely wider on average. The bootstrap resamples
day-blocks around a single leaky multiplier fit once on the full matched
test set; the interval is conditional on that multiplier and quantifies
sampling variability in the width difference, not variability in the
multiplier itself. TimesFM-3's real advantage over
one-parameter rescaling is therefore narrower than a CRPS-only reading
suggests, and it sits on a different axis than CRPS: on CRPS, a trivial
multiplier already beats it, zero-shot; the width/coverage trade-off is
where TimesFM-3 holds an edge that neither a leakage-free nor a
leakage-optimistic scalar rescaling reproduces. Section~\ref{sec:discussion}
returns to what this split result means for claims that trivial baselines
match foundation models generally \citep{zhang_2025_context_parroting}.

A separate experiment supplies TimesFM-3 with ensemble spread as a
second, simultaneous covariate (Table~\ref{tab:full-comparison}'s
two-covariate row). CRPS worsens to \CrpsTsfmTwoCov{} K and coverage
moves closer to the 0.80-nominal target (\CoverageTsfmTwoCov{} against
\CoverageTsfmFull{}) at a substantially wider interval
(\WidthTsfmTwoCov{} K) -- a shift along the calibration-sharpness
trade-off, not a clean improvement, discussed fully in
Section~\ref{sec:discussion-covariate-check}.

\subsection{The training-data-size breakpoint curve}
\label{sec:results-breakpoint}

\begin{figure}[htbp]
\centering
\includegraphics[width=0.95\linewidth]{figures/f1_breakpoint_curve}
\caption{CRPS, coverage@80\%, and interval width against training data
size (cases, log scale, with calendar-day equivalent on the secondary
axis). Horizontal lines are N-independent zero-shot references; curves are
trained methods refit at each N.}
\label{fig:f1}
\end{figure}

Figure~\ref{fig:f1} shows the paper's headline result: EMOS crosses
TimesFM-3's CRPS at \BpCrpsPooledCases{} cases pooled
(\BpCrpsPooledDays{} calendar days) and \BpCrpsLocalCases{} cases locally
(\BpCrpsLocalDays{} days). That raw crossing is a point estimate, not
itself statistically significant: testing significance at every
low-\(N\) point on this curve (\(k=1,2,3,5,7,9\), Section~\ref{sec:results-significance})
shows the advantage is significant from \BpCrpsPooledSigK{} cases
(\BpCrpsPooledSigDays{} calendar days) onward for pooled EMOS and from
\BpCrpsLocalSigK{} cases (\BpCrpsLocalSigDays{} days) for local EMOS --
under day-blocking, our primary block definition; the station-blocked
\(t\)-test does not reach significance anywhere in the tested
\(k=1\)--9 range for either variant, while the station-blocked Wilcoxon
test does, from \(k=\BpCrpsPooledStationWilcoxonFirstSigK{}\) onward --
the same test-statistic divergence already reported at \SigTestNCases{}
cases (Section~\ref{sec:results-significance}), now shown to persist
across the whole low-\(k\) range. Both numbers matter and neither
replaces the other: the raw crossing (\BpCrpsPooledDays{} days) is where
the two methods' CRPS point estimates change order; the day-blocked
significant crossing (\BpCrpsPooledSigDays{} days) is where that
difference is established rather than plausibly noise. DRN needs
\BpCrpsDrnCases{} cases (\BpCrpsDrnDays{} days) -- under the
feature-constrained setup this paper trains it in, not DRN's
originally-intended full feature set (Section~\ref{sec:limitations}).
The variance-inflation
baseline (variant a, CRPS-optimal, refit per training size the same way
EMOS is) is already better at the smallest tested N: at
\InflASmallestNDays{} days it already reaches \InflASmallestNCrps{} K,
below TimesFM-3's \CrpsTsfmFull{} K, and stays there throughout the
tested range -- the same qualitative pattern Table~\ref{tab:full-comparison}
shows at full data (Section~\ref{sec:results-full}). Coverage tells a
different story (ratios below computed from the unrounded breakpoint
estimates, not the values as displayed): EMOS pooled crosses TimesFM-3's
coverage at \BpCoveragePooledCases{} cases (\BpCoveragePooledDays{} days)
-- \CoverageCrpsRatioPooled{}x later than its CRPS crossing -- and EMOS
local not until \BpCoverageLocalCases{} cases (\BpCoverageLocalDays{}
days), \CoverageCrpsRatioLocal{}x later than its own CRPS crossing; DRN crosses
at \BpCoverageDrnCases{} cases (\BpCoverageDrnDays{} days). Interval width
has no crossing at all (Section~\ref{sec:breakpoint-def}): EMOS pooled's interval
starts narrower than TimesFM-3's and crosses to wider somewhere between
90 and 365 training days, while EMOS local stays narrower than TimesFM-3
across the entire tested range even as it reaches comparable coverage --
a sharpness advantage a pooled fit does not share.

\subsection{Lead-time resolution}
\label{sec:results-leadtime}

\begin{figure}[htbp]
\centering
\includegraphics[width=0.95\linewidth]{figures/f2_lead_time_resolved}
\caption{CRPS and coverage@80\% by lead time (21 steps, 0--120~h).}
\label{fig:f2}
\end{figure}

\begin{figure}[htbp]
\centering
\includegraphics[width=0.95\linewidth]{figures/f5_breakpoint_by_lead_group}
\caption{Training-data-size breakpoint, computed separately within each
lead-time bucket, for CRPS and coverage@80\%.}
\label{fig:f5}
\end{figure}

Pooling over all 21 lead times hides a reversal (Figure~\ref{fig:f2}):
TimesFM-3's CRPS advantage is real but short-lived, with a first sign flip
at \FirstSignFlipHours{}~h and a durable crossover at
\DurableCrossoverHours{}~h, after which TimesFM-3 remains the worse method
at every longer tested lead time. Recomputing the breakpoint curve
separately within three lead-time buckets (0--24~h, 24--72~h, 72--120~h;
Figure~\ref{fig:f5}) sharpens this into the paper's most lead-dependent
result: at 72--120~h, both EMOS pooled and EMOS local already beat
TimesFM-3 on CRPS at the smallest tested training size (a single case);
at 0--24~h, pooled EMOS does not cross TimesFM-3's CRPS anywhere in the
tested range, from 1 to 4180 cases. Coverage crossings follow the same
ordering (earliest at long lead, latest at short lead) but at
substantially larger training sizes than CRPS in every bucket.

\subsection{Sharpness, calibration, and PIT diagnostics}
\label{sec:results-sharpness}

\begin{figure}[htbp]
\centering
\includegraphics[width=0.7\linewidth]{figures/f3_sharpness_calibration_plane}
\caption{Interval width against coverage@80\% for every method; trained
methods shown as trajectories across training size, zero-shot methods as
single points.}
\label{fig:f3}
\end{figure}

\begin{figure}[htbp]
\centering
\includegraphics[width=0.95\linewidth]{figures/f4_pit_histograms}
\caption{PIT histograms at full training data. PIT uniformity is assessed
descriptively; no formal uniformity test valid under serial dependence is
applied.}
\label{fig:f4}
\end{figure}

Figure~\ref{fig:f3} makes Section~\ref{sec:results-full}'s width/coverage
trade-off visible directly: EMOS local's trajectory sits up and to the
left of every other trained method's, reaching comparable coverage to
TimesFM-3 with a narrower interval throughout. The raw ensemble's PIT
histogram (Figure~\ref{fig:f4}) shows the classic severe-underdispersion
signature -- roughly a third of observations fall in each of the two
outermost bins -- consistent with its \CoverageRawFull{} empirical
coverage against a 0.80 nominal target; TimesFM-3 and both EMOS variants
show comparatively flat histograms, read descriptively per
Section~\ref{sec:methods}.

Table~\ref{tab:full-comparison}'s twCRPS column follows the same ranking
as CRPS itself for every method (e.g.\ EMOS local's
\TwcrpsEmosLocalFull{} K remains the lowest, TimesFM-3's
\TwcrpsTsfmFull{} K remains among the highest) -- this checks the upper
end of the quantile grid this study fixed for every method, not whether
TimesFM-3 genuinely lacks tail behavior beyond p90
(Section~\ref{sec:metrics}).

\subsection{Statistical significance}
\label{sec:results-significance}

At the pooled headline comparison (EMOS pooled vs.\ TimesFM-3 at
\SigTestNCases{} cases, \SigTestNDays{} days), the CRPS advantage is robust to block
definition: station-blocked and day-blocked tests agree in direction, and
the day-blocked test -- which additionally absorbs cross-station synoptic
dependence the station block cannot -- gives \(p=\SigCrpsDayT\) (paired
\(t\)) and \(p=\SigCrpsDayW\) (Wilcoxon), both far below the station-blocked
test's already-significant Wilcoxon result (\(p=\SigCrpsStationW\)) and its
non-significant \(t\)-test (\(p=\SigCrpsStationT\)). The coverage advantage
does not clear the same bar: station-blocked (\(t\): \(p=\SigCoverageStationT\);
Wilcoxon: \(p=\SigCoverageStationW\)) and day-blocked
(\(t\): \(p=\SigCoverageDayT\); Wilcoxon: \(p=\SigCoverageDayW\)) tests
disagree with each other on both significance and, at the day block,
direction of evidence between the two test statistics. We do not report
EMOS as coverage-superior to TimesFM-3 at this training size: the honest
summary is that the CRPS result is significant under every tested block
definition and test statistic, while the coverage result is not.

The breakpoint curve's own raw crossing (Section~\ref{sec:results-breakpoint}:
\BpCrpsPooledCases{} cases pooled, \BpCrpsLocalCases{} cases locally) is a
different, smaller training size than \SigTestNCases{} cases, and was not
itself significance-tested until now. Repeating the day- and
station-blocked tests at every point on the low-\(N\) breakpoint grid
(\(k=1,2,3,5,7,9\)) shows the CRPS advantage is not significant at the
raw crossing: the earliest \(k\) at which both the day-blocked \(t\)-test
and Wilcoxon test (day-blocking is our primary block definition) agree
the advantage is significant is
\BpCrpsPooledSigK{} cases (\BpCrpsPooledSigDays{} days) for pooled EMOS
(day-blocked \(t\): \(p=\BpCrpsPooledSigKDayT\); Wilcoxon:
\(p=\BpCrpsPooledSigKDayW\)) and \BpCrpsLocalSigK{} cases
(\BpCrpsLocalSigDays{} days) for local EMOS (day-blocked \(t\):
\(p=\BpCrpsLocalSigKDayT\); Wilcoxon: \(p=\BpCrpsLocalSigKDayW\)) --
meaningfully later than the raw crossing, though still a small absolute
training size. The station-blocked paired \(t\)-test is not significant
at this same \(k\) for either variant (pooled: \(p=\BpCrpsPooledSigKStationT\);
local: \(p=\BpCrpsLocalSigKStationT\)), while the station-blocked Wilcoxon
test is (pooled: \(p=\BpCrpsPooledSigKStationW\); local:
\(p=\BpCrpsLocalSigKStationW\)) -- the same station-block
\(t\)-vs-Wilcoxon divergence already seen at \SigTestNCases{} cases above
holds throughout the tested range, not just at that one point; the
station-blocked \(t\)-test is not significant at any \(k\) we tested from
1 to 9 cases for either EMOS variant.

Repeating both block definitions separately within each lead-time bucket,
using full-training-data EMOS pooled against TimesFM-3, resolves this
divergence differently per bucket rather than uniformly. At 0--24~h, the
CRPS gap is significant in TimesFM-3's favor (day-blocked \(t\):
\(p=\LeadSigCrpsBucketShortDayT\); station:
\(p=\LeadSigCrpsBucketShortStationT\)). The coverage gap in the same
bucket is significant in EMOS's favor instead (day-blocked \(t\):
\(p=\LeadSigCoverageBucketShortDayT\); station:
\(p=\LeadSigCoverageBucketShortStationT\)) -- EMOS is significantly
\emph{closer to nominal coverage} even in the one bucket where TimesFM-3
wins on CRPS. At 24--72~h, station- and day-blocked
CRPS tests disagree on significance (day: \(p=\LeadSigCrpsBucketMidDayT\);
station: \(p=\LeadSigCrpsBucketMidStationT\)), with the day-blocked test --
our primary block definition -- favoring EMOS. At 72--120~h, both block
definitions agree EMOS's CRPS advantage is significant
(\(p=\LeadSigCrpsBucketLongDayT\) day-blocked,
\(p=\LeadSigCrpsBucketLongStationT\) station-blocked). Across all three
buckets, EMOS's coverage advantage over TimesFM-3 is significant under
day-blocking (\(p\) between \LeadSigCoverageBucketShortDayT{} and
\LeadSigCoverageBucketLongDayT{}). The pattern across every test in this
section is consistent in direction and variable in magnitude and
significance, not uniform: we report each result at the resolution it was
actually tested, rather than a single pooled significance claim standing
in for all of them.

\section{Discussion}
\label{sec:discussion}

\subsection{Why a few days is enough}
\label{sec:discussion-fewdays}

The paper's central practical result is that EMOS pooled crosses
TimesFM-3's CRPS at \BpCrpsPooledCases{} cases (\BpCrpsPooledDays{}
calendar days, Section~\ref{sec:results-breakpoint}) -- but that raw
crossing is a point estimate, not yet a statistically defensible claim.
Testing significance across the low-\(N\) range shows the earliest
training size at which the advantage is significant under both a
day-blocked \(t\)-test and a Wilcoxon test is \BpCrpsPooledSigK{} cases
(\BpCrpsPooledSigDays{} days, Section~\ref{sec:results-significance}) --
day-blocking being our primary block definition; the station-blocked
\(t\)-test never reaches significance in the tested \(k=1\)--9 range,
while the station-blocked Wilcoxon test does from
\(k=\BpCrpsPooledStationWilcoxonFirstSigK{}\) onward, the same
divergence already reported at \SigTestNCases{} cases now shown to
persist throughout -- still a small absolute training size, and still
well inside the regime
\citet{baran_2016_combining} and \citet{gebetsberger_2018_estimation}
already identify as usable for station-adaptive EMOS, and that the
time-series and lead-time-continuous EMOS variants of
\citet{jobst_2024_ts_emos} and \citet{wessel_2024_leadtime_continuous}
target directly -- but nearly twice the raw crossing's calendar-day
figure, not the same number. The practical implication for a newly
commissioned station is direct, stated at the number the evidence
actually supports: roughly \BpCrpsPooledSigDays{} days of paired data,
not months, are enough to stop relying on a zero-shot foundation model's
raw CRPS with statistical confidence; the raw crossing at
\BpCrpsPooledDays{} days is the earliest the two methods' CRPS point
estimates change order, not the earliest that difference is established
rather than noise.

This result is deliberately scoped to CRPS alone. Coverage tells a
different story at substantially more data -- \CoverageCrpsRatioPooled{}x
later for pooled EMOS, \CoverageCrpsRatioLocal{}x later for local EMOS
(ratios from the unrounded breakpoint estimates, not the values as
displayed; Section~\ref{sec:results-breakpoint}: \BpCoveragePooledCases{} cases
pooled, \BpCoverageLocalCases{} cases locally), and interval width has no
single crossing point at all -- EMOS pooled starts sharper than TimesFM-3
and crosses to wider as it is given more data, while EMOS local stays
sharper throughout the tested range (Section~\ref{sec:breakpoint-def}'s
three-outcome distinction exists precisely so this kind of metric-specific
divergence is reported as what it is, not flattened into one number). "A
few days is enough" is therefore true for CRPS specifically, not for every
metric this paper reports -- Sections~\ref{sec:discussion-trivial}
and~\ref{sec:discussion-mechanism} address what CRPS alone does and does
not establish about why.

The question this paper asks -- how much training data a trained method
needs to replace a zero-shot foundation model -- has a deployment-cost
dimension this paper has not systematically measured: TimesFM-3 requires
a GPU forward pass through a 330-million-parameter pretrained model for
every inference call, while fitting EMOS end-to-end on the full training
archive took \EmosPooledFitSeconds{}~s (pooled) or
\EmosLocalFitSeconds{}~s (local) of CPU time in this study's own runs, and
inference afterward is a closed-form Gaussian evaluation, not a further
optimization. Once the training-data threshold in this paper is crossed,
the practical choice is not a single-axis trade-off between accuracy and
convenience: the trained method is both more accurate and cheaper to run.

\subsection{Positioning against trivial baselines}
\label{sec:discussion-trivial}

Section~\ref{sec:results-full}'s full-data comparison (Table~\ref{tab:full-comparison})
shows a result that splits by metric in a way worth stating plainly rather
than softening: on CRPS, a one-parameter variance-inflation rescaling of
the raw ensemble already beats TimesFM-3, and the fixed-climatological
variant (b, \CrpsInflBFull{} K vs.\ TimesFM-3's \CrpsTsfmFull{} K) does so
with no training data at all -- squarely the pattern
\citet{zhang_2025_context_parroting} describes for trivial baselines
against zero-shot foundation models generally, now confirmed in station
postprocessing specifically. A comparable dynamic -- a pretrained model
matching, but not decisively beating, a much simpler well-tuned
alternative -- is reported in a different domain by
\citet{molinaro_2026_universal_downscaling}.

On the calibration-sharpness axis, the evidence points the other way,
and this finding is now settled rather than conditional: the
coverage-matched variant (c), its multiplier fit on training data alone
(never on test-set coverage, Section~\ref{sec:methods}), needs
\WidthInflCFull{} K to reach as close to TimesFM-3's own coverage as a
training-fit multiplier allows -- substantially wider than TimesFM-3's
\WidthTsfmFull{} K. Variant (c) is leakage-free by design, which means it
is not an upper bound on what a trivial one-parameter rescaling could
achieve at TimesFM-3's coverage target -- only a lower bound on how well
a \emph{legitimately deployable} one can do. The leakage-optimistic
counterpart -- a multiplier fit directly on the test set's own coverage,
answering "what is the best any trivial rescaling could ever do here,
even with an unfair advantage" -- needs \InflLeakyWidth{} K, still wider
than TimesFM-3's \WidthTsfmFull{} K, at matched coverage: the leaky
variant's root-finder converges to \InflLeakyCoverage{}, identical to
TimesFM-3's \CoverageTsfmFull{} target to within floating-point
precision, so the width comparison is not confounded by a residual
coverage gap the way variant (c)'s is. This margin is not a bare point
estimate: a day-blocked bootstrap on the per-instance width difference
gives \InflLeakyWidthDiffPoint{} K (95\% CI
[\InflLeakyWidthDiffCiLow{}, \InflLeakyWidthDiffCiHigh{}]) -- entirely
positive, so the difference is statistically distinguishable from zero,
not merely numerically positive on average. The interval is conditional
on the leakage-optimistic multiplier fit once on the full test set, not
re-fit within each bootstrap resample: it quantifies sampling variability
in the width difference at that fixed multiplier, not variability in the
multiplier itself. This margin is also not explained by EMOS's own
training/test ensemble-size mismatch: subsampling the test side down to
the training side's member count changes EMOS's own interval width by no
more than \EnsSubsampleWidthDiffLocal{} K (Section~\ref{sec:limitations}),
under one-tenth of this margin. TimesFM-3's
calibration-sharpness advantage is therefore real: it survives even this
unrealistically favorable trivial baseline, with uncertainty quantified
(subject to that conditioning), so \citet{zhang_2025_context_parroting}'s
critique, while correct on CRPS (above), does not extend to this axis.

The honest summary of Section~\ref{sec:results-full} is: TimesFM-3 loses
to a trivial rescaling on CRPS, unconditionally; it beats even a
leakage-optimistic trivial rescaling on sharpness at matched coverage,
so the calibration-sharpness advantage is not a leakage-free-vs.-leaky
comparison artifact.

\subsection{The short-lead advantage: information asymmetry or architecture}
\label{sec:discussion-mechanism}

\textbf{This section remains provisional; the architectural-advantage
hypothesis has not been confirmed.} Section~\ref{sec:covariates}
proposed a single mechanism for TimesFM-3's short-lead CRPS advantage
(Section~\ref{sec:results-leadtime}: first sign flip at
\FirstSignFlipHours{}~h, durable crossover at
\DurableCrossoverHours{}~h): access to two channels EMOS and DRN were not
given in this study's configuration (Table~\ref{tab:info-channels}) --
the station's own recent observation history, and ensemble spread. These
are two separate questions with two separate tests; one has now been
run, with a limited, inconclusive result, and one is still outstanding.

\textbf{Observation history (tested; does not support the
information-asymmetry explanation, but does not settle the question
either).} A persistence-augmented EMOS at 0--24~h tests whether adding
the station's own last observation as a predictor closes the CRPS gap
toward TimesFM-3. It does not, substantially: only
\PersistenceGapClosedPct{}\% of the gap closes (persistence weight
\PersistenceWeight{}, matched subset, \NPersistenceMatchedInstances{}
instances), confirmed after ruling out lead-0 contamination as a
confound (Section~\ref{sec:limitations}: exact-match and trend checks
both clean). This result does not support the information-asymmetry
explanation as stated -- the simplest way of giving EMOS the missing
observation-history channel does not close the gap. It does not, on its
own, establish that the advantage is architectural instead: the test has
three real limitations that keep it from distinguishing the two
hypotheses cleanly, not just describing them as a formality.

\begin{itemize}
\item \emph{Lead-pooled fit.} EMOS pools all \NLeadTimes{} lead times
into one fit (Section~\ref{sec:methods-list}); the persistence term is
fit once, lead-blind, inside that pooled model, while TimesFM-3's actual
covariate and target information vary naturally with lead time. The two
methods are not given the same structural opportunity to use
lead-dependent information -- this is not an equal comparison.
\item \emph{One linear term, one lag.} The persistence predictor is a
single linear term on the single most recent observation. TimesFM-3
receives a context window of 40 or more past observations, processed by
a sequence model that can learn arbitrarily nonlinear, multi-lag
structure. Adding one linear lag does not give EMOS "comparable
information" to a 40-observation context in any meaningful sense --
only a much smaller, cruder approximation to it.
\item \emph{Metric-dependent.} In the same 0--24~h bucket, TimesFM-3
wins on CRPS but EMOS is significantly \emph{closer to nominal coverage}
(day-blocked \(t\): \(p=\LeadSigCoverageBucketShortDayT\), EMOS's favor,
Section~\ref{sec:results-significance}). Whatever "architectural
advantage" this test might support holds for CRPS specifically, not
uniformly across every metric in the same lead-time bucket.
\end{itemize}

Given these three limitations, the honest reading of this test is
narrower than "architecture, not information asymmetry": it shows that
the simplest information-asymmetry fix is insufficient, not that
architecture is confirmed as the explanation. A test able to distinguish
the two hypotheses would need a lead-resolved EMOS fit (not pooled) with
a multi-lag persistence predictor (not one linear term on one lag) --
this paper does not run that test, and does not claim the weaker test
run here settles the question.

\textbf{Ensemble spread / calibration (resolved: does not support the
asymmetry framing).} The other half of the proposed asymmetry -- that
TimesFM-3's calibration deficit relative to EMOS reflects a missing
spread channel, not an architectural limit -- has now been tested
directly. The model's covariate channel does accept spread as a second,
simultaneous input, and the model is genuinely sensitive to that
channel's real informational content, not merely its presence (a
preliminary test confirms both, Section~\ref{sec:covariates},
Section~\ref{sec:limitations}), but supplying it this way did not close
the calibration deficit: CRPS and interval width both worsen, and
coverage moves only marginally closer to nominal at a substantially
wider interval (Section~\ref{sec:discussion-covariate-check} reports the
full result). This does not confirm the asymmetry framing for this half
of the mechanism question, and the observation-history result above does
not settle it either: the two channels test genuinely different things,
and a finding for one is not evidence for the other.

\subsection{The two-covariate evaluation: does spread change calibration?}
\label{sec:discussion-covariate-check}

Section~\ref{sec:covariates}'s preliminary test (T1) established that
TimesFM-3's covariate channel accepts ensemble spread as a second,
simultaneous input alongside the ensemble mean, and that the model's
output is genuinely sensitive to that channel's real informational
content. This section reports the full-scale evaluation that
preliminary result motivated: supplying both the ensemble mean and
ensemble spread as simultaneous past-future covariates, on the same
\NMatchedInstances{} matched instances used for every other TimesFM-3
number in this paper, and comparing the result against the mean-only
configuration used everywhere else.

Supplying spread did not improve TimesFM-3's CRPS: it worsens from
\CrpsTsfmFull{} K to \CrpsTsfmTwoCov{} K, a net negative of
\CrpsTwoCovDiff{} K. No uncertainty interval has been computed for this
difference; its magnitude is comparable to the ensemble-size sensitivity
measured directly in Section~\ref{sec:limitations}
(\EnsSubsampleCrpsDiffPooled{}--\EnsSubsampleCrpsDiffLocal{} K), so it
should be read as a real but not dramatically large effect, not a
precisely bounded one. Coverage moves closer to the 0.80-nominal target,
from \CoverageTsfmFull{} to \CoverageTsfmTwoCov{}: in practice, the
two-covariate configuration is calibrated at this target, where the
mean-only configuration under-covers. Interval width also increases,
from \WidthTsfmFull{} K to \WidthTsfmTwoCov{} K (\WidthTwoCovDiff{} K
wider) -- but by this paper's own standard of judging sharpness jointly
with calibration (Section~\ref{sec:methods},
\citealp{gneiting_2007_calibration_sharpness}), these two widths are not
directly comparable on the same axis: they are measured at different
coverage levels, and a wider interval is the expected cost of a
configuration that is closer to properly calibrated, not on its own
evidence of a worse model. This reading has independent support: fitting
the leakage-optimistic multiplier (Section~\ref{sec:discussion-trivial})
at this same \CoverageTsfmTwoCov{} target, rather than TimesFM-3's
original \CoverageTsfmFull{}, still needs \LeakyWidthAtTwoCovCoverage{}
K -- \LeakyWidthAtTwoCovCoverageDiff{} K wider than the two-covariate
configuration's own \WidthTsfmTwoCov{} K at the identical coverage.
Finding 2's width advantage therefore holds at two independent coverage
points, not only at TimesFM-3's original one.

Supplying ensemble spread as a second covariate therefore worsens
TimesFM-3's CRPS and moves it along the calibration-sharpness trade-off
-- closer to nominal coverage, at a wider interval -- rather than
cleanly closing the calibration deficit: this is a shift along that
trade-off, not a clean improvement or a clean failure. This result does
not confirm the information-asymmetry framing for the calibration half
of the mechanism question (Section~\ref{sec:covariates}): TimesFM-3 was
not simply prevented by an input-design choice from recovering EMOS-like
calibration once that choice was corrected. But it does not confirm
architecture as the explanation instead, either: the model accepted the
second input and was demonstrably sensitive to it
(Section~\ref{sec:limitations}), and did shift its output along the
calibration-sharpness trade-off when given it -- it just did not convert
that shift into a net improvement on CRPS, or into a clean win on the
paper's joint sharpness-calibration standard, within this study's
zero-shot, frozen-weights setup. Whether a fine-tuned or differently
prompted configuration would do better remains untested and outside this
paper's scope; so does whether a different input scale or normalization
convention for the spread channel would change this result
(Section~\ref{sec:limitations}).

\subsection{Implementation pitfall: covariate wiring}
\label{sec:pitfall}

\citet{kreusel_2026_covariates_key} reports that covariate handling, not
architecture or scale, is the dominant factor in time series foundation
model performance on real tasks. This project's own development history is
a direct, first-hand instance of that finding rather than a citation of
it: an earlier version of TimesFM-3's covariate pipeline
(Section~\ref{sec:covariates}) supplied an incoherent covariate channel --
one that ran without error and produced plausible-looking quantile output
-- and measurably degraded TimesFM-3's apparent skill before the bug was
found and corrected. The pipeline did not crash; it silently produced a
worse model. The practical warning for anyone implementing a
covariate-aware time series foundation model pipeline is that "it runs and
returns sensible-looking numbers" is not evidence the covariate is wired
correctly -- only a dedicated correctness check on the covariate channel
itself (not just on the model's output shape or its overall CRPS) would
have caught this before it silently degraded results. This experience is
the direct reason this project later required a distinguishing
correctness check -- not just "does the call succeed," but "is the
output genuinely sensitive to this covariate's real informational
content, not merely its presence" -- before trusting any further
covariate-channel change (Section~\ref{sec:discussion-covariate-check}'s
ablation test is that check, run precisely because of the lesson learned
here).

\subsection{The value of seasonal coverage in small training sets}
\label{sec:discussion-seasonal}

Section~\ref{sec:data}'s two sampling arms answer different questions by
construction: the contiguous arm's small-\(k\) subsets see less seasonal
variety than a full archive would, while the random arm's uniform draw
gives artificial full-season coverage even at small \(k\). At
\SeasonalCompareNDays{} days, EMOS pooled's coverage is already visibly
different between the two arms -- \CoverageRandomSeasonalN{} under the
random arm against \CoverageContigSeasonalN{} under the contiguous arm --
while CRPS is nearly identical (\CrpsRandomSeasonalN{} K vs.\
\CrpsContigSeasonalN{} K). This is consistent with
Section~\ref{sec:discussion-fewdays}'s CRPS-crosses-fast,
coverage-crosses-slow pattern having a seasonal-coverage component: a
model can already predict the mean well from a handful of contiguous days,
but calibrating the predictive interval correctly needs to have seen a
representative spread of conditions, not just recent ones. Operationally,
this means the realistic deployment scenario -- a newly commissioned
station with only contiguous recent history, the contiguous arm's premise
-- reaches good CRPS quickly but should not be assumed to be well
calibrated at the same training size a full-season random sample would be.
\citet{baran_2026_sharpness_penalty} targets the same sharpness axis on
the same benchmark; the overlap is in the axis studied, not in this
specific contiguous-vs-random comparison, which does not appear in that
work.

\subsection{Limitations}
\label{sec:limitations}

\textbf{Covariate channel design; tested, but does not resolve the
mechanism question.} TimesFM-3's covariate channel was given only the
ensemble mean in this study's main comparison
(Section~\ref{sec:covariates}); whether supplying ensemble spread too
would change calibration was tested directly
(Section~\ref{sec:discussion-covariate-check}). The channel does accept
a second, simultaneous input, and the model is genuinely sensitive to
its real informational content, not merely its presence -- but CRPS
worsens when spread is supplied this way, and while coverage moves
closer to nominal (a genuine calibration shift, not merely a wider
interval for its own sake), the interval widens substantially to get
there, so the net result is a shift along the calibration-sharpness
trade-off, not a clean win (Section~\ref{sec:discussion-covariate-check}).
This result does not confirm the information-asymmetry framing (spread
was not simply withheld from a channel that would otherwise have used it
straightforwardly beneficially), but it does not confirm architecture as
the explanation either: the channel accepted the input and the model was
demonstrably sensitive to it, and did shift its behavior in response to
it, without that shift converting into a net improvement. Whether a
fine-tuned or differently prompted configuration would use spread more
successfully is untested and outside this paper's scope. We report this
as an unresolved mechanism question rather than resolving it either way
beyond what this evidence supports.

\textbf{Spread channel scale and normalization.} Ensemble spread was
supplied to TimesFM-3's second covariate channel as a standard
deviation, in the same physical units (Kelvin) as the ensemble mean but
roughly two orders of magnitude smaller in raw value (typically 1--3~K
vs.\ 270--290~K) -- a real scale disparity between the two channels as
supplied. This is not, on its own, an unaddressed confound: TimesFM-3's
own preprocessing z-normalizes each covariate channel independently,
computing a running mean and standard deviation per channel rather than
sharing statistics across channels (confirmed directly from the
installed \texttt{timesfm3} package's source, not inferred), so the raw
magnitude difference between the mean and spread channels does not by
itself explain why supplying spread failed to improve performance. The
spread channel was supplied at a single scale and as a single channel,
per-channel normalization is confirmed rather than merely assumed, and
the result above should be read as "this specific input design does not
help," not as "TimesFM-3 cannot use ensemble spread" more broadly --
untested alternatives include a differently transformed spread channel
(e.g.\ log-spread) or a channel expressing spread relative to the mean.

\textbf{Lead-time pooling.} EMOS and DRN each fit a single parameter set
pooled across all \NLeadTimes{} lead times (Section~\ref{sec:methods-list}),
while TimesFM-3 and the raw ensemble vary naturally by lead time. A
lead-blind fit is necessarily too wide at short lead (where true
uncertainty is small) and too narrow at long lead (where it is large)
relative to a lead-time-aware fit, which means the CRPS crossings reported
in Section~\ref{sec:results-breakpoint} are, if anything, understated for
EMOS and DRN, not artifacts of this pooling. It is a genuine limitation
for the sharpness comparison specifically, since a lead-blind interval's
width behavior is not directly comparable to TimesFM-3's per-lead-varying
output. For DRN in particular, this feature-constrained setup is not the
architecture's originally-intended one (\citealp{rasp_2018_drn} additionally
use auxiliary NWP fields and lead time as inputs): where DRN underperforms
in this paper, the correct reading is that this architecture, under this
feature-constrained setup, was outperformed -- not that the architecture
itself is weak.

\textbf{Persistence test run on a lead-pooled EMOS fit.} The
persistence-augmented EMOS test (Section~\ref{sec:discussion-mechanism})
inherits this same lead-pooling limitation directly: the persistence
predictor is fit once, lead-blind, inside a model already disadvantaged
by pooling across all 21 lead times against a naturally lead-varying
TimesFM-3. This reduces the test's power to distinguish an
information-asymmetry explanation from an architectural one -- a
lead-resolved EMOS fit with a multi-lag persistence predictor would be a
meaningfully stronger test, and this paper does not run one.

\textbf{Target-variable specificity.} The breakpoint measured here is
for \(2\,\mathrm{m}\) temperature, where EMOS's Gaussian predictive
family is close to ideally specified. For targets requiring non-Gaussian
or target-specific reformulations (Section~\ref{sec:intro}: e.g.\
\citealp{scheuerer_2013_precip_emos}'s precipitation-specific EMOS), the
trained method's small-sample advantage would not necessarily appear as
early.

\textbf{Tail truncation in twCRPS.} The tail-weighted CRPS defined in
Section~\ref{sec:metrics} uses the upper two levels of a fixed
\NQuantileLevels{}-level quantile grid (0.8, 0.9) as its tail weighting.
This tests the upper end of the quantile grid this study fixed for every
method, not whether TimesFM-3 genuinely lacks tail behavior beyond p90 --
the grid itself has no quantile levels past 0.9 for any method, so no
comparison in this paper can distinguish "no tail" from "a tail this
grid cannot see."

\textbf{PIT diagnostics are descriptive only.}
\citet{broecker_2018_serial_dependence} and
\citet{broecker_2020_stratified_rank} show that formal rank-uniformity
hypothesis tests require correction for serial dependence, which this
study's PIT values genuinely exhibit (station~$\times$~time~$\times$~lead
panel structure, Section~\ref{sec:significance}); no such test is applied
to the PIT histograms in Section~\ref{sec:results-sharpness}, which are
read descriptively for this reason.
\citet{dirkson_2025_misdiagnosing_reliability} further shows that coverage
and rank-based reliability diagnostics can both be misled by the same
covariance pathology -- we do not claim the descriptive PIT reporting in
this paper is immune to that critique; it is a genuine, unresolved
limitation, not one worked around by reporting PIT histograms instead of
a formal test.

\textbf{Random-arm seed count.} The random sampling arm's mean and
standard deviation (Section~\ref{sec:data}) are estimated from
\DataSizeSweepSeeds{} seeds. This is enough to see the arm's qualitative
pattern relative to the contiguous arm, but the standard deviation itself
is estimated noisily from so few draws -- it should be read as an
indicator of spread, not a precise quantity, and Figure~\ref{fig:f1}'s
random-arm error bars inherit this same noise.

\textbf{Ensemble size mismatch.} EMOS's training-side ensemble variance is
estimated from \TrainEnsembleMembers{} members against
\TestEnsembleMembers{} on the test side (Section~\ref{sec:data}); an
11-member variance estimate is intrinsically noisier than a 51-member one
of the same true spread, so EMOS is trained against a noisier variance
estimate than the one used to evaluate it. This is no longer an
unbounded concern: subsampling 11 of the test side's 51 members (10
draws, recomputing the ensemble mean/variance from just those 11 and
re-evaluating the already-fitted full-\(N\) EMOS models) changes pooled
EMOS's own CRPS by \EnsSubsampleCrpsDiffPooled{} K, its interval width by
\EnsSubsampleWidthDiffPooled{} K, and its coverage@80\% by
\EnsSubsampleCoverageDiffPooled{} (local EMOS: \EnsSubsampleCrpsDiffLocal{} K,
\EnsSubsampleWidthDiffLocal{} K, \EnsSubsampleCoverageDiffLocal{}
respectively) -- both width differences are under one-tenth of
Section~\ref{sec:discussion-trivial}'s \InflLeakyWidthDiffPoint{} K
leakage-optimistic-width margin. The ensemble-size mismatch is real and
measurable, but it is too small to be a material contributor to that
finding.

\textbf{Unverified chronological assumption.} The contiguous sampling
arm orders cases by descending \((\mathrm{year\_idx}, \mathrm{time\_idx})\),
assuming ascending \(\mathrm{year\_idx}\) is chronologically more recent
(Section~\ref{sec:data}). Only \(\mathrm{time\_idx}\)'s ordering was
independently verified against real decoded dates; \(\mathrm{year\_idx}\)'s
polarity was not. This bounds its impact to training sizes beyond the
209-date single-cycle boundary (\(k>209\), i.e.\ N greater than roughly
728 calendar days) -- below that boundary, every contiguous-arm result in
this paper is unaffected regardless of which polarity is correct.

\textbf{Training-data contamination risk.} This study does not verify that
TimesFM-3's pretraining corpus excludes the 2017--2018 test period used
here. Data contamination in time series foundation model evaluation is an
active, documented concern
(\citealp{li_2026_tsfmaudit,chiazagomekpere_2026_leakage_aware,balloccu_2024_leak_cheat_repeat}),
and benchmark-methodology work
(\citealp{aksu_2024_gift_eval,meyer_2025_rethinking_eval}) argues evaluation
protocols in this area have not caught up with the risk. We cannot rule
this out for TimesFM-3 on EUPPBench; a demonstrated case of the same risk
materializing in a different domain is given by
\citet{pan_2026_contamination_electricity}. If TimesFM-3's apparent
zero-shot skill benefits from contamination, its true zero-shot CRPS would
be worse than reported here, which would only strengthen this paper's
conclusion that trained methods overtake it quickly -- it does not
threaten the paper's central claim, but it is a real, unaddressed
limitation of the TimesFM-3 numbers taken in isolation.

\textbf{Lead-0 contamination, checked and ruled out.}
Section~\ref{sec:discussion-mechanism} depends on the 0--24~h bucket,
which includes \(\mathrm{step\_hours}=0\) -- the analysis time, where the
"forecast" could in principle be tightly coupled to the observation it
is verified against for every method, not only TimesFM-3. Two checks
rule this out: the exact-match fraction between the ensemble-mean
forecast and the observation at lead 0 is \LeadZeroExactMatchFraction{}
(no higher than the \LeadZeroMaxOtherExactMatchFraction{} seen at any other
lead, both consistent with floating-point noise, not literal copying);
and lead 0's residual standard deviation (\LeadZeroActualStd{} K) is in
line with what a linear trend fit to leads \(\geq 6\)~h predicts at
\(\mathrm{step\_hours}=0\) (\LeadZeroPredictedStd{} K, ratio
\LeadZeroStdRatio{}) -- if anything very slightly above the trend, not a
discontinuous tightening below it. Every lead-time-resolved result in
this paper that includes lead 0 (Section~\ref{sec:results-leadtime}) can
be read without this particular caveat.

\section{Conclusion}
\label{sec:conclusion}

At full training data, every trained method -- EMOS pooled, EMOS local,
and DRN -- beats TimesFM-3's CRPS, and a one-parameter variance-inflation
rescaling of the raw ensemble does too, with a fixed-climatological
variant achieving this using no training data at all
(Section~\ref{sec:results-full}). This advantage appears at very small
training sizes: EMOS pooled's raw CRPS crossing is at
\BpCrpsPooledCases{} cases (\BpCrpsPooledDays{} calendar days), and that
crossing becomes statistically significant, under day-blocking (our
primary block definition), from \BpCrpsPooledSigK{} cases
(\BpCrpsPooledSigDays{} days) onward for both EMOS variants -- the
station-blocked \(t\)-test does not reach significance anywhere in the
tested \(k=1\)--9 range, while the station-blocked Wilcoxon test does
from \(k=\BpCrpsPooledStationWilcoxonFirstSigK{}\) onward
(Section~\ref{sec:results-significance}). For a newly commissioned
station, this is the practical answer to this paper's title question on
CRPS: roughly \BpCrpsPooledSigDays{} days of paired ensemble/observation
history, not months, are enough to reach a statistically defensible CRPS
advantage over a frozen, zero-shot foundation model, and does so at a
lower inference cost too, on the qualitative comparison in
Section~\ref{sec:discussion-fewdays} (a real, measured EMOS fit time
against an unmeasured but qualitatively larger GPU inference cost for
TimesFM-3, not a quantified comparison on both sides). Coverage crosses
far more slowly
than CRPS -- \CoverageCrpsRatioPooled{}x later for pooled EMOS,
\CoverageCrpsRatioLocal{}x later for local EMOS -- so this practical
answer holds for CRPS specifically, not for calibration.

On the calibration-sharpness axis, TimesFM-3 holds a real advantage over
a trivial rescaling, and this is settled, not conditional: even a
leakage-optimistic multiplier fit directly on test-set coverage -- an
upper bound no legitimately deployable trivial rescaling could exceed --
needs \InflLeakyWidth{} K, still wider than TimesFM-3's \WidthTsfmFull{}
K at matched coverage, by a margin of \InflLeakyWidthDiffPoint{} K (95\%
CI [\InflLeakyWidthDiffCiLow{}, \InflLeakyWidthDiffCiHigh{}],
day-blocked bootstrap, Section~\ref{sec:discussion-trivial}) that
excludes zero entirely. This advantage is not an
artifact of comparing a leakage-free multiplier against a
leakage-optimistic ceiling; it survives even the most favorable trivial
baseline this paper could construct.

TimesFM-3's short-lead CRPS advantage reverses within the tested
lead-time range -- first sign flip at \FirstSignFlipHours{}~h, durable
reversal by \DurableCrossoverHours{}~h
(Section~\ref{sec:results-leadtime}) -- and the mechanism behind that
short-lead advantage has two independent parts, both now tested
separately. The observation-history part does not support an
information-asymmetry explanation: a persistence-augmented EMOS at
0--24~h closes only \PersistenceGapClosedPct{}\% of the gap toward
TimesFM-3. This result is not treated as confirming an architectural
explanation instead, since the test itself has three limitations that
keep it from cleanly distinguishing the two hypotheses -- a lead-pooled
EMOS fit, a single linear persistence term against TimesFM-3's much
richer context window, and a metric-dependent result (EMOS is
significantly closer to nominal coverage in the same bucket where
TimesFM-3 wins on CRPS) (Section~\ref{sec:discussion-mechanism}). The
ensemble-spread/calibration part has also now been tested directly:
TimesFM-3's covariate channel does accept a second, simultaneous input,
and the model is genuinely sensitive to that channel's real
informational content, not merely its presence
(Section~\ref{sec:covariates}, Section~\ref{sec:limitations}) -- but
supplying ensemble spread this way does not cleanly close the
calibration deficit. CRPS worsens from \CrpsTsfmFull{} K to
\CrpsTsfmTwoCov{} K; coverage moves closer to nominal, from
\CoverageTsfmFull{} to \CoverageTsfmTwoCov{} -- a genuine calibration
shift, not merely a wider interval -- but interval width also increases
substantially, from \WidthTsfmFull{} K to \WidthTsfmTwoCov{} K, the
expected cost of that shift under this paper's own standard of judging
sharpness jointly with calibration, not independent evidence of a worse
model (Section~\ref{sec:discussion-covariate-check}). The net effect is
a shift along the calibration-sharpness trade-off together with a worse
CRPS, not a clean improvement or a clean failure. Neither half of the
mechanism therefore supports the information-asymmetry framing -- but
neither confirms architecture as the explanation either: the channel
accepted the second input and was demonstrably sensitive to it, and did
shift its behavior in response, yet that shift did not convert into a
net improvement within this study's zero-shot, frozen-weights setup. The
short-lead advantage's mechanism remains an open interpretive question,
with architecture one candidate explanation among others, not a
confirmed one.

Two questions that were open earlier in this project are now both
resolved. EMOS's training/test ensemble-size mismatch is bounded and too
small to be a material contributor to the calibration-sharpness
comparison (Section~\ref{sec:limitations}). And whether ensemble spread
could be supplied to TimesFM-3 as a second covariate, and whether doing
so would change its calibration, has been tested directly
(Section~\ref{sec:discussion-covariate-check}): it can be supplied, but
doing so does not improve calibration in a way that supports the
information-asymmetry framing (Section~\ref{sec:discussion-mechanism}).
Neither resolution changes this paper's two central findings: on CRPS, a
frozen, zero-shot time series foundation model is
not competitive with a few days of station-specific training data, for
either a trained statistical postprocessing method or a trivial
one-parameter rescaling of the raw ensemble (though the exact training
size at which this becomes statistically significant, rather than a raw
point-estimate crossing, depends on block definition,
Section~\ref{sec:results-significance}); and on calibration-sharpness,
TimesFM-3's advantage over even a leakage-optimistic trivial rescaling is
real. The evidence for both is for one time series foundation model,
TimesFM-3; whether it generalizes to other pretrained time series models
(e.g.\ Chronos-2, Moirai-2) is untested. Both findings are also specific
to the target variable studied here, \(2\,\mathrm{m}\) temperature, where
EMOS's Gaussian predictive family is close to ideally specified; for
targets requiring non-Gaussian or target-specific reformulations, a
trained method's small-sample advantage should not be assumed to appear
as early.


\section*{Code and Data Availability}
\addcontentsline{toc}{section}{Code and Data Availability}

Code for this study is publicly available at
\href{https://github.com/elendmire/zero-shot-postprocessing}{github.com/elendmire/zero-shot-postprocessing}.
%% D30 (panel review, 2026-09-07): PLACEHOLDER sentence below. Code is
%% still changing (pending SSH runs, see docs/pending_ssh_runs.md) --
%% do not mint the Zenodo archival release until the analysis is frozen
%% for submission, and replace the sentence below with the real DOI
%% once it exists. Never invent a DOI here.
A versioned, archival release with its own DOI will be minted via Zenodo
before journal submission and cited here once it exists.
EUPPBench, version \EuppbenchVersion{}, is described in
\citet{demaeyer_2023_euppbench} and distributed via Zenodo (DOI:
10.5281/zenodo.7429236, verified by direct DOI resolution and Zenodo
record inspection, 2026-09-07); this study uses its own defined
train/test split without modification (Section~\ref{sec:data}). TimesFM
3.0 weights (checkpoint \texttt{google/timesfm-3.0-pytorch}) are obtained
from the model card at
\url{https://huggingface.co/google/timesfm-3.0-pytorch} (Section~\ref{sec:covariates}),
distributed under the TimesFM Non-Commercial License v1.0 (verified
against the model card's own license field, 2026-09-07); this study's use
is academic and non-commercial, consistent with that license.

\section*{Funding}
\addcontentsline{toc}{section}{Funding}

This research received no specific grant from any funding agency in the
public, commercial, or not-for-profit sectors.

\section*{CRediT Author Contributions}
\addcontentsline{toc}{section}{CRediT Author Contributions}

\"{O}mer Faruk Avc{\i}: Conceptualization, Methodology, Software,
Validation, Formal analysis, Investigation, Data curation, Writing --
original draft, Writing -- review \& editing, Visualization. Sole author;
all roles.

\section*{Competing Interests}
\addcontentsline{toc}{section}{Competing Interests}

The author declares no competing interests.

\section*{Acknowledgements}
\addcontentsline{toc}{section}{Acknowledgements}

Computations in this study were performed on \.{I}stanbul Technical
University's Altay high-performance computing cluster. Claude (Anthropic,
Claude Sonnet 5) was used to assist with code implementation; the author
is responsible for all scientific content.


\bibliographystyle{plainnat}
\bibliography{refs}

\end{document}
